% Vietnamese translation for OI Public Library
% Source-PDF: IOI中国国家候选队论文/2002/李澎煦.pdf
\documentclass[11pt]{article}
\usepackage{oipl-vn}

\title{Thuật toán giao nửa mặt phẳng và ứng dụng}
\author{Li Pengxu}
\date{}

\begin{document}
\maketitle

\origtitle{Algorithms and applications of half-plane intersection}
\sourcepath{IOI China candidate team papers / 2002 / Li Pengxu.pdf}

\section*{Tổng quan}

Bài viết giới thiệu các khái niệm cơ bản của nửa mặt phẳng và giao nửa mặt
phẳng, rồi trình bày hai cách tính giao nửa mặt phẳng: thuật toán trực tuyến
dựa trên việc lần lượt cắt miền hiện có, và thuật toán chia để trị dựa trên
việc hợp nhất giao của hai tập nửa mặt phẳng. Phần ứng dụng cho thấy nhiều
bài toán hình học trong thi lập trình có thể được chuyển thành bài toán tìm
giao của các nửa mặt phẳng hoặc xét tính không rỗng của nghiệm một hệ bất
đẳng thức tuyến tính hai biến.

\paragraph{Từ khóa.}
Nửa mặt phẳng; giao nửa mặt phẳng; đa giác lồi; chia để trị; Voronoi.

\tableofcontents

\section{Khái niệm cơ bản}

\term{Nửa mặt phẳng} là một đường thẳng trên mặt phẳng cùng với phần nằm ở
một phía của đường thẳng đó. Trong hệ tọa độ Descartes, một nửa mặt phẳng có
thể được xác định bởi bất đẳng thức
\[
  ax+by+c\ge 0.
\]

Trong một miền bị chặn, khi tính toán thực tế ta có thể giả sử có một biên đủ
lớn, một nửa mặt phẳng hoặc giao của các nửa mặt phẳng là một miền đa giác
lồi.

Giao của \(n\) nửa mặt phẳng
\[
  H_1\cap H_2\cap \cdots \cap H_n
\]
là một đa giác lồi có nhiều nhất \(n\) cạnh.

\section{Thuật toán}

\subsection{Thuật toán trực tuyến cho giao nửa mặt phẳng}

\begin{verbatim}
procedure intersection of half-planes
Input:
  n half-planes H1,H2,...,Hn, corresponding to inequalities
  ai*x+bi*y+ci >= 0, i=1,2,3,...,n
Output:
  H1 cap H2 cap ... cap Hn

  Initialize region A as the whole plane.
  For i = 1,2,...,n:
    Cut A by the line ai*x+bi*y+ci = 0.
    Keep the part satisfying ai*x+bi*y+ci >= 0.
  Output A.
\end{verbatim}

Thuật toán này có độ phức tạp thời gian \(O(n^2)\), đồng thời có ưu điểm của
một thuật toán trực tuyến: sau mỗi nửa mặt phẳng mới, ta có thể cập nhật ngay
miền giao hiện tại.

\subsection{Thuật toán chia để trị cho giao nửa mặt phẳng}

Giả sử ta có thể hợp nhất giao của \(m\) nửa mặt phẳng và giao của \(n\) nửa
mặt phẳng trong thời gian \(O(m+n)\). Khi đó có thể dùng một thuật toán chia
để trị \(O(n\log n)\) để tìm giao của các nửa mặt phẳng.

\begin{verbatim}
Procedure intersection of half-plane (D&C)
Input:
  n half-planes H1,H2,...,Hn, corresponding to inequalities
  ai*x+bi*y+ci >= 0, i=1,2,3,...,n
Output:
  H1 cap H2 cap ... cap Hn

  Split H1,...,Hn into two sets of approximately equal size.
  Recursively compute the intersection in each subproblem.
  Merge the two convex polygon regions to form
  H1 cap H2 cap ... cap Hn.
\end{verbatim}

Vì vậy, điểm mấu chốt của bài toán là làm thế nào để tìm giao của hai đa giác
lồi trong thời gian \(O(m+n)\).

Ta cắt hai đa giác lồi theo các đường song song với trục \(y\). Trong thời
gian \(O(m+n)\), hai đa giác có thể được chia thành nhiều nhất \(O(m+n)\)
miền hình thang. Giao của hai hình thang tương ứng có thể được xử lý trong
thời gian \(O(1)\).

Để thuận tiện cho thao tác, ta dùng một cách đặc biệt để biểu diễn đa giác
lồi. Có thể thấy các đỉnh ở biên trên và các đỉnh ở biên dưới của một đa
giác lồi lần lượt tạo thành hai dãy có tọa độ \(x\) tăng dần. Lưu riêng hai
dãy đỉnh này dưới dạng hai danh sách liên kết sẽ cho ta biên trên và biên
dưới xác định miền đa giác lồi.

\paragraph{Thuật toán.}
\begin{verbatim}
procedure intersection of convex polygon
Input:
  two convex polygon regions A and B
Output:
  C = A cap B

1. Classify the vertices of the two polygons by x-coordinate, obtaining
   the sequence x_i, i=1,...,p.
2. Initialize C as the empty region.
3. Process {x_1}.
4. Process the regions (x_i, x_{i+1}], i=1,...,p-1, in order.
   4.1 Compute the trapezoids, possibly degenerate, cut from the two
       polygons in this region. Call them ABCD and A'B'C'D'.
   4.2 Compute the intersection points
       AB cap A'B', AB cap C'D', and CD cap A'B'.
       Sort all existing points by x-coordinate, delete duplicates, and
       append them to the upper boundary of C. Compute the lower boundary
       of C in the same way.
   4.3 Compute the right boundary of this region:
       EF = BC cap B'C'. Append E and F to the upper and lower boundaries
       of C, respectively.
5. Output C.
\end{verbatim}

Ở bước 1, vì tọa độ \(x\) trên biên trên và biên dưới của \(A\) và \(B\) đều
đã có thứ tự, ta có thể dùng thao tác trộn tương tự merge sort. Độ phức tạp
là \(O(m+n)\).

Ở bước 4, các miền được quét theo thứ tự tọa độ \(x\) tăng dần, nên con trỏ
trên mỗi biên luôn chỉ dịch sang phải trong toàn bộ quá trình. Mỗi đỉnh của
hai đa giác được quét nhiều nhất một lần. Độ phức tạp là \(O(m+n)\).

Do đó, độ phức tạp thời gian của toàn bộ thuật toán là \(O(m+n)\).

\section{Ứng dụng}

\subsection{Bài toán 1: Hotter and Colder (Waterloo local contest)}

\paragraph{Lược thuật đề bài.}
Hai người \(A\) và \(B\) chơi một trò chơi trên bàn cờ \(10\times 10\). Người
\(A\) xác định một điểm \(P\), còn \(B\) di chuyển một lần trong mỗi lượt.
Mỗi lần, \(A\) nói cho \(B\) biết vị trí hiện tại của \(B\) gần \(P\) hơn
(\textit{Hot}) hay xa \(P\) hơn (\textit{Cold}) so với trước đó. Đề gốc còn
phải xét trường hợp khoảng cách không đổi.

Sau mỗi câu trả lời của \(A\), hãy xác định diện tích miền mà điểm \(P\) còn
có thể nằm trong đó.

\paragraph{Phân tích.}
Giả sử \(B\) di chuyển từ \(C(x_1,y_1)\) đến \(D(x_2,y_2)\), và \(A\) trả lời
\textit{Hot}. Khi đó, đối với lượt này, vị trí \(P(x,y)\) phải thỏa
\(|CP|>|DP|\), tức là
\[
  (2x_2-2x_1)x+(2y_2-2y_1)y+x_1^2+y_1^2-x_2^2-y_2^2>0.
\]
Tương tự, câu trả lời \textit{Cold} tương ứng với một bất đẳng thức khác.

Ban đầu, miền có thể là \([0,10]\times[0,10]\). Sau mỗi lượt, ta lấy giao
của miền hiện tại với nửa mặt phẳng tương ứng với bất đẳng thức mới, rồi xuất
diện tích của giao đó.

\subsection{Bài toán 2: Milk (OOPC1)}

\paragraph{Lược thuật đề bài.}
SRbGa có một miếng bánh mì hình đa giác lồi \(n\) cạnh, và một chậu sữa có
diện tích đủ lớn nhưng độ sâu chỉ bằng \(h\). Anh ta muốn chỉ nhúng bánh
\(k\) lần, mỗi lần đều bảo đảm miếng bánh vuông góc với đáy chậu, sao cho
diện tích phần bánh dính sữa là lớn nhất.

\paragraph{Phân tích.}
Do quy mô bài này không lớn, ta xét dùng tìm kiếm theo chiều sâu.

Nhúng theo mỗi cạnh tương ứng với một nửa mặt phẳng của phần còn lại. Với một
cách chọn \(k\) cạnh \(E_1,\ldots,E_k\) để nhúng, phần bánh còn lại chính là
giao của \(k\) nửa mặt phẳng này với đa giác ban đầu. Ta xét \(C(n,k)\) cách
nhúng và chọn cách làm cho diện tích còn lại nhỏ nhất.

Bài toán 1 dùng một số nửa mặt phẳng để lấy giao tuần tự, và sau mỗi lần đều
phải xuất diện tích. Rõ ràng thuật toán trực tuyến là phù hợp.

Với bài toán 2, nếu dùng thuật toán trực tuyến thì độ phức tạp là
\(O(C(n,k)\,n)\), đồng thời tiện cắt nhánh trong quá trình tìm kiếm. Nếu dùng
thuật toán chia để trị ngoại tuyến, độ phức tạp là
\[
  O\bigl(C(n,k)(n+k\log k)\bigr).
\]

\subsection{Bài toán 3: Video (CTSC98)}

\paragraph{Lược thuật đề bài.}
Cho một đa giác \(P\), không nhất thiết lồi. Hỏi trong \(P\) có tồn tại một
điểm \(Q\) sao cho từ \(Q\) có thể quan sát toàn bộ miền đa giác hay không.

\paragraph{Phân tích.}
Giả sử các điểm biên của đa giác được cho theo thứ tự ngược chiều kim đồng
hồ:
\[
  V_0,V_1,V_2,\ldots,V_n,\qquad V_0=V_n.
\]
Một điểm \(Q_i\) có thể quan sát cạnh \(V_iV_{i+1}\) nhất định phải thỏa
\[
  \overrightarrow{Q_iV_i}\times
  \overrightarrow{Q_iV_{i+1}}\ge 0,\qquad i=0,\ldots,n-1.
\]

Hơn nữa, một điểm có thể quan sát tất cả các cạnh thì nhất định có thể quan
sát toàn bộ miền đa giác.

Nếu dùng tọa độ để tính tích có hướng, mỗi ràng buộc tương ứng với một bất
đẳng thức bậc nhất hai biến, cũng tức là một nửa mặt phẳng. Vì vậy bài toán
được chuyển thành việc xét giao của \(n\) nửa mặt phẳng này có rỗng hay
không.

\subsection{Bài toán 4: Triathlon (NEERC2000)}

\paragraph{Lược thuật đề bài.}
Có \(n\) vận động viên tham gia cuộc thi ba môn phối hợp. Thứ hạng được xác
định theo tổng thời gian mỗi vận động viên dùng trên ba chặng. Biết vận tốc
của vận động viên \(i\) trong ba môn lần lượt là \(U_i,V_i,W_i\).

Hỏi đối với vận động viên \(i\), có thể sắp xếp thích hợp độ dài của ba chặng
thi, trong đó độ dài mỗi chặng đều không được bằng \(0\), để bảo đảm vận động
viên đó thắng hay không.

\paragraph{Phân tích.}
Giả sử độ dài ba chặng lần lượt là \(x,y,z\). Điều kiện cần và đủ để vận động
viên \(i\) thắng là, với mọi \(j\ne i\),
\[
  \frac{x}{U_i}+\frac{y}{V_i}+\frac{z}{W_i}
  <
  \frac{x}{U_j}+\frac{y}{V_j}+\frac{z}{W_j}.
\]

Đây là một hệ bất đẳng thức thuần nhất ba biến. Vì \(z>0\), ta có thể chia
hai vế của mỗi bất đẳng thức cho \(z\), rồi đặt \(X=x/z\), \(Y=y/z\). Khi đó
ta nhận được
\[
  \left(\frac{1}{U_j}-\frac{1}{U_i}\right)X
  +\left(\frac{1}{V_j}-\frac{1}{V_i}\right)Y
  +\left(\frac{1}{W_j}-\frac{1}{W_i}\right)>0,
  \qquad j\ne i.
\]

Bài toán được chuyển thành việc tìm giao của \(n-1\) nửa mặt phẳng tương ứng
với các bất đẳng thức này, rồi kiểm tra diện tích giao có lớn hơn \(0\) hay
không, tức là loại bỏ các trường hợp tập rỗng, một điểm, hoặc một đoạn thẳng.

Bài toán 3 và bài toán 4 cuối cùng đều được chuyển thành bài toán tồn tại
nghiệm của một hệ bất đẳng thức hai biến. Có thể dùng thuật toán chia để trị
để giải tương đối hiệu quả.

\subsection{Bài toán 5: Run away (CERC99)}

\paragraph{Lược thuật đề bài.}
Trong một hình chữ nhật \(R\) có \(n\) điểm \(P_1,\ldots,P_n\). Hãy tìm một
điểm \(Q\in R\) sao cho
\[
  \min_i |QP_i|
\]
là lớn nhất.

\paragraph{Phân tích.}
Chia \(R\) thành \(n\) miền \(Q_1,\ldots,Q_n\), trong đó \(Q_i\) là tập các
điểm gần \(P_i\) hơn so với mọi điểm còn lại:
\[
  Q_i=\{\,Q\mid |QP_i|\le |QP_j|,\ i\ne j\,\}\cap R.
\]

Miền \(Q_i\) có thể được tìm bằng giao của các nửa mặt phẳng nằm về phía
\(P_i\) đối với đường trung trực của \(P_iP_j\). Vì vậy \(Q_i\) là một đa
giác lồi.

Trong \(Q_i\), điểm xa \(P_i\) nhất chỉ có thể xuất hiện ở một đỉnh của
\(Q_i\). Do đó, chỉ cần tìm điểm xa nhất trong các đỉnh ấy.

Tìm giao nửa mặt phẳng bằng thuật toán chia để trị có độ phức tạp
\(O(n\log n)\). Đa giác tương ứng với \(P_i\) có nhiều nhất \(O(n)\) đỉnh,
nên việc tìm điểm xa nhất trong \(Q_i\) có độ phức tạp \(O(n)\).

Độ phức tạp tổng cộng là \(O(n^2\log n)\).

Thực ra, \(n\) miền đa giác \(Q_1,\ldots,Q_n\) được định nghĩa theo phương
pháp trên tạo thành một biểu đồ Voronoi. Biểu đồ Voronoi là một đối tượng
hình học trong hình học tính toán, có tầm quan trọng chỉ sau bao lồi, và có
ứng dụng rất rộng rãi. Phương pháp dùng giao nửa mặt phẳng để tìm biểu đồ
Voronoi không phải phương pháp tối ưu. Nhiều thuật toán như chia để trị hoặc
quét mặt phẳng đều có thể đạt độ phức tạp \(O(n\log n)\), và đó là tối ưu.
Tuy nhiên, các thuật toán ấy quá phức tạp và không thuộc phạm vi thảo luận
của bài viết này.

\section{Tài liệu tham khảo}

\begin{enumerate}
  \item F. P. Preparata và M. I. Shamos, \textit{Nhập môn hình học tính toán},
  bản in lần thứ nhất, tháng 11 năm 1990.
  \item Zhou Peide, \textit{Hình học tính toán: phân tích và thiết kế thuật
  toán}, bản in lần thứ nhất, tháng 3 năm 2000.
\end{enumerate}

\appendix
\section{Văn bản bản chiếu trích xuất}

Phần sau chuyển sang tiếng Việt phần văn bản trích xuất được từ các bản
chiếu nằm sau bài viết trong tệp PDF.

\subsection*{Trang tiêu đề}

\begin{center}
\Large Thuật toán giao nửa mặt phẳng và ứng dụng

\normalsize Trường Trung học số 4 Bắc Kinh

Li Pengxu
\end{center}

\subsection*{Khái niệm cơ bản}

\begin{itemize}
  \item Nửa mặt phẳng: đường thẳng trên mặt phẳng và phần nằm ở một phía của
  nó.
  \item Nửa mặt phẳng có thể được xác định bởi bất đẳng thức
  \(ax+by+c\ge 0\).
  \item Trong một miền bị chặn, nửa mặt phẳng hoặc giao của các nửa mặt phẳng
  là một miền đa giác lồi.
  \item Giao của \(n\) nửa mặt phẳng là một đa giác lồi có nhiều nhất \(n\)
  cạnh.
\end{itemize}

\subsection*{Thuật toán trực tuyến cho giao nửa mặt phẳng}

\begin{verbatim}
procedure intersection of half-planes
Input:
  n half-planes H1,H2,...,Hn
Output:
  H1 cap H2 cap ... cap Hn

  Initialize region A as the whole plane.
  Cut A successively with the lines ai*x+bi*y+ci=0,
  i=1,2,...,n, keeping the part satisfying ai*x+bi*y+ci>=0.
  Output A.

Complexity O(n*n); online algorithm.
\end{verbatim}

\subsection*{Thuật toán chia để trị cho giao nửa mặt phẳng}

Giả sử có thể hợp nhất giao của \(m\) nửa mặt phẳng và giao của \(n\) nửa
mặt phẳng trong thời gian \(O(m+n)\). Khi đó có một thuật toán chia để trị
\(O(n\log n)\) để tìm giao nửa mặt phẳng.

\begin{verbatim}
Procedure intersection of half-plane (D&C)
Input:
  n half-planes H1,H2,...,Hn
Output:
  H1 cap H2 cap ... cap Hn

  Split H1,...,Hn into two sets of approximately equal size.
  Recursively compute the intersection in each subproblem.
  Merge the two convex polygon regions to form
  H1 cap H2 cap ... cap Hn.
\end{verbatim}

Điểm mấu chốt là làm thế nào để tìm giao của hai đa giác lồi trong thời gian
\(O(m+n)\).

\begin{itemize}
  \item Cắt hai đa giác lồi theo các đỉnh thành nhiều nhất \(O(m+n)\) hình
  thang song song với trục \(y\).
  \item Giao của mỗi cặp hình thang có thể được xử lý trong thời gian
  \(O(1)\).
\end{itemize}

\subsection*{Cách mô tả đa giác lồi}

\begin{itemize}
  \item Các đỉnh ở phía trên và phía dưới của một đa giác lồi lần lượt tạo
  thành hai dãy có tọa độ \(x\) tăng dần.
  \item Lưu hai dãy đỉnh này dưới dạng hai danh sách liên kết, ta nhận được
  biên trên và biên dưới xác định miền đa giác lồi.
\end{itemize}

\subsection*{Thuật toán giao đa giác lồi, phần 1}

\begin{verbatim}
procedure intersection of convex polygon
Input:
  two convex polygon regions A and B
Output:
  C = A cap B

1. Classify the vertex x-coordinates of the two polygons, obtaining
   the sequence x_i, i=1,...,p.
2. Initialize C as empty.
3. Process {x_1}.
4. Process the regions (x_i,x_{i+1}], i=1,...,p-1, in order.
5. Output C.
\end{verbatim}

\subsection*{Thuật toán giao đa giác lồi, phần 2}

\begin{verbatim}
4. Process the regions (x_i,x_{i+1}], i=1,...,p-1, in order.
   4.1 Compute the trapezoids, possibly degenerate, cut from the two
       polygons in this region: ABCD and A'B'C'D'.
   4.2 Compute AB cap A'B', AB cap C'D', and CD cap A'B'.
       Sort the existing points by x-coordinate, delete duplicates,
       and append them to the upper boundary of C.
   4.3 Compute the lower boundary of C in the same way.
   4.4 Compute the right boundary of this region, as a segment
       intersection: EF = BC cap B'C'. Append E and F to the upper
       and lower boundaries of C, respectively.
\end{verbatim}

\subsection*{Độ phức tạp của thuật toán}

\begin{itemize}
  \item Bước 1: vì tọa độ \(x\) trên biên trên và biên dưới của \(A\), \(B\)
  đều đã có thứ tự, có thể dùng thao tác trộn. Độ phức tạp \(O(m+n)\).
  \item Bước 4: vì các miền được quét theo thứ tự \(x\) tăng dần, con trỏ
  trên mỗi biên luôn dịch sang phải trong toàn bộ quá trình. Mỗi đỉnh của
  hai đa giác được quét nhiều nhất một lần. Độ phức tạp \(O(m+n)\).
  \item Độ phức tạp thời gian của toàn bộ thuật toán là \(O(m+n)\).
\end{itemize}

\subsection*{Bài toán 1: Hotter and Colder}

Waterloo local contest.

Hai người \(A\) và \(B\) chơi trên bàn cờ \(10\times 10\). Người \(A\) xác
định một điểm \(P\), người \(B\) di chuyển một lần trong mỗi lượt. Mỗi lần,
\(A\) nói cho \(B\) biết vị trí hiện tại gần \(P\) hơn (\textit{Hot}) hay xa
\(P\) hơn (\textit{Cold}). Đề gốc còn phải xét trường hợp khoảng cách không
đổi. Sau mỗi câu trả lời, hãy xác định diện tích miền mà \(P\) có thể nằm
trong đó.

\subsection*{Phân tích bài toán 1}

\begin{itemize}
  \item Giả sử \(B\) đi từ \(C(x_1,y_1)\) đến \(D(x_2,y_2)\), và \(A\) trả
  lời \textit{Hot}. Khi đó \(P(x,y)\) thỏa \(|CP|>|DP|\), tức là
  \[
    2(x_2-x_1)x+2(y_2-y_1)y+x_1^2+y_1^2-x_2^2-y_2^2>0.
  \]
  \item Tương tự, \textit{Cold} tương ứng với một bất đẳng thức khác.
  \item Ban đầu, miền có thể là \([0,10]\times[0,10]\). Sau mỗi lượt, lấy
  giao miền hiện tại với nửa mặt phẳng tương ứng, rồi xuất diện tích giao.
\end{itemize}

\subsection*{Bài toán 2: Nice Milk}

OOPC1.

SRbGa có một miếng bánh mì hình đa giác lồi \(n\) cạnh, và một chậu sữa có
diện tích đủ lớn nhưng độ sâu chỉ bằng \(h\). Anh ta muốn chỉ nhúng \(k\) lần,
mỗi lần đều bảo đảm bánh vuông góc với đáy chậu, sao cho diện tích phần bánh
dính sữa là lớn nhất.

\subsection*{Phân tích bài toán 2}

\begin{itemize}
  \item Do quy mô bài này không lớn, xét dùng tìm kiếm theo chiều sâu.
  \item Nhúng theo mỗi cạnh tương ứng với một nửa mặt phẳng của phần còn lại.
  Với một cách nhúng \(k\) cạnh \(E_1,\ldots,E_k\), phần còn lại tương ứng
  với giao của \(k\) nửa mặt phẳng ấy và đa giác ban đầu.
  \item Xét \(C(n,k)\) cách nhúng và chọn cách có diện tích còn lại nhỏ nhất.
\end{itemize}

\subsection*{Tiểu kết}

\begin{itemize}
  \item Bài toán 1 dùng một số nửa mặt phẳng để lấy giao tuần tự, và mỗi lần
  đều phải xuất diện tích. Thuật toán trực tuyến là phù hợp.
  \item Với bài toán 2, nếu dùng thuật toán trực tuyến, độ phức tạp là
  \(O(C(n,k)n)\), đồng thời tiện cắt nhánh trong quá trình tìm kiếm. Nếu dùng
  thuật toán chia để trị ngoại tuyến, độ phức tạp là
  \(O(C(n,k)(n+k\log k))\).
\end{itemize}

\subsection*{Bài toán 3: Video}

CTSC98.

Cho một đa giác \(P\), không nhất thiết lồi. Hỏi trong \(P\) có tồn tại điểm
\(Q\) sao cho từ \(Q\) quan sát được toàn bộ miền đa giác hay không.

\subsection*{Phân tích bài toán 3}

\begin{itemize}
  \item Nếu các đỉnh đa giác được cho theo thứ tự ngược chiều kim đồng hồ
  \(V_0,V_1,V_2,\ldots,V_n\), \(V_0=V_n\), thì điểm \(Q_i\) có thể quan sát
  cạnh \(V_iV_{i+1}\) nhất định thỏa
  \[
    \overrightarrow{Q_iV_i}\times\overrightarrow{Q_iV_{i+1}}\ge 0,
    \qquad i=0,\ldots,n-1.
  \]
  \item Điểm quan sát được mọi cạnh nhất định quan sát được toàn bộ miền đa
  giác.
  \item Dùng tọa độ để tính tích có hướng, mỗi ràng buộc tương ứng với một
  bất đẳng thức bậc nhất hai biến, tức là một nửa mặt phẳng. Bài toán trở
  thành xét giao của \(n\) nửa mặt phẳng có rỗng hay không.
\end{itemize}

\subsection*{Bài toán 4: Triathlon}

NEERC2000.

Có \(n\) vận động viên tham gia ba môn phối hợp. Thứ hạng được sắp theo tổng
thời gian dùng trong ba chặng. Biết vận tốc của mỗi vận động viên ở ba môn
là \(U_i,V_i,W_i\). Hỏi đối với vận động viên \(i\), có thể sắp xếp thích
hợp độ dài ba chặng, nhưng độ dài mỗi chặng đều không được bằng \(0\), để bảo
đảm người đó thắng hay không.

\subsection*{Phân tích bài toán 4}

\begin{itemize}
  \item Giả sử độ dài ba chặng là \(x,y,z\). Điều kiện cần và đủ để vận động
  viên \(i\) thắng là
  \[
    \frac{x}{u_i}+\frac{y}{v_i}+\frac{z}{w_i}
    <
    \frac{x}{u_j}+\frac{y}{v_j}+\frac{z}{w_j}.
  \]
  \item Đây là hệ bất đẳng thức thuần nhất ba biến. Vì \(z>0\), chia hai vế
  cho \(z\), rồi đặt \(X=x/z\), \(Y=y/z\), ta được
  \[
    \left(\frac{1}{u_j}-\frac{1}{u_i}\right)X
    +\left(\frac{1}{v_j}-\frac{1}{v_i}\right)Y
    +\left(\frac{1}{w_j}-\frac{1}{w_i}\right)>0.
  \]
  \item Bài toán trở thành tìm giao của \(n-1\) nửa mặt phẳng tương ứng và
  kiểm tra diện tích có lớn hơn \(0\) hay không, tức là loại bỏ tập rỗng,
  điểm và đoạn thẳng.
\end{itemize}

\subsection*{Tiểu kết}

\begin{itemize}
  \item Bài toán 3 và bài toán 4 cuối cùng đều trở thành bài toán tồn tại
  nghiệm của một hệ bất đẳng thức hai biến. Có thể dùng thuật toán chia để
  trị của giao nửa mặt phẳng để giải hiệu quả.
  \item Hai bài toán hơi khác nhau: một bên dùng biên nghiêm ngặt, một bên
  dùng biên không nghiêm ngặt. Nói cách khác, cách xử lý biên của đa giác là
  khác nhau. Với bất đẳng thức có dấu \(\ge 0\), cần xét các trường hợp suy
  biến thành điểm hoặc đoạn thẳng, nên phức tạp hơn một chút.
\end{itemize}

\subsection*{Bài toán 5: Run away}

CERC99.

Trong một hình chữ nhật \(R\) có \(n\) điểm \(P_1,\ldots,P_n\). Hãy tìm một
điểm \(Q\in R\) sao cho \(\min_i |QP_i|\) là lớn nhất.

\subsection*{Phân tích bài toán 5, phần 1}

\begin{itemize}
  \item Chia \(R\) thành \(n\) miền \(Q_1,\ldots,Q_n\), trong đó \(Q_i\) là
  tập các điểm trong \(R\) có khoảng cách đến \(P_i\) không lớn hơn khoảng
  cách đến các điểm khác:
  \[
    Q_i=\{\,Q\mid |QP_i|\le |QP_j|,\ i\ne j\,\}\cap R.
  \]
  \item \(Q_i\) có thể được tìm bằng giao của các nửa mặt phẳng nằm về phía
  \(P_i\) đối với đường trung trực của \(P_iP_j\). Vì vậy \(Q_i\) là một đa
  giác lồi.
  \item Trong \(Q_i\), điểm xa \(P_i\) nhất chỉ có thể là một đỉnh của
  \(Q_i\). Tìm điểm xa nhất trong các đỉnh đó là đủ.
\end{itemize}

\subsection*{Phân tích bài toán 5, phần 2}

\begin{itemize}
  \item Dùng thuật toán chia để trị cho giao nửa mặt phẳng, độ phức tạp cho
  mỗi điểm là \(O(n\log n)\).
  \item Đa giác tương ứng với \(P_i\) có nhiều nhất \(O(n)\) đỉnh, nên tìm
  điểm xa nhất trong \(Q_i\) có độ phức tạp \(O(n)\).
  \item Độ phức tạp tổng cộng là \(O(n^2\log n)\).
\end{itemize}

\subsection*{Biểu đồ Voronoi}

\begin{itemize}
  \item Thực ra, \(n\) miền đa giác \(Q_1,\ldots,Q_n\) được định nghĩa ở trên
  tạo thành một biểu đồ Voronoi.
  \item Biểu đồ Voronoi là một đối tượng hình học trong hình học tính toán,
  có tầm quan trọng chỉ sau bao lồi, và có ứng dụng rất rộng rãi.
  \item Thuật toán dùng giao nửa mặt phẳng để tìm biểu đồ Voronoi không phải
  là tối ưu.
  \item Chia để trị, quét mặt phẳng và nhiều thuật toán khác đều đạt được độ
  phức tạp \(O(n\log n)\), đó mới là tối ưu. Nhưng các thuật toán ấy quá phức
  tạp và không thuộc phạm vi thảo luận của bài viết này.
\end{itemize}

\subsection*{Dàn ý}

\begin{itemize}
  \item Khái niệm cơ bản.
  \item Thuật toán.
  \item Thuật toán trực tuyến cho giao nửa mặt phẳng.
  \item Thuật toán chia để trị cho giao nửa mặt phẳng.
  \item Ứng dụng.
  \item Bài toán 1: Hotter and Colder.
  \item Bài toán 2: Milk.
  \item Tiểu kết.
  \item Bài toán 3: Video.
  \item Bài toán 4: Triathlon.
  \item Tiểu kết.
  \item Bài toán 5: Run away.
  \item Biểu đồ Voronoi.
\end{itemize}

\end{document}
